\appendix
\addtocontents{toc}{\protect\setcounter{tocdepth}{1}}
\crefalias{section}{appendix}
\crefname{appendix}{Appendix}{Appendices}

\section{A combinatorial model for simple--simple derived Hom}
\label{app:simple-simple-assembly}

In this appendix, we prove
\Cref{prop:simple-simple-combinatorial-model}.
Throughout this appendix, $D(\kk)$ denotes the derived category of
cochain complexes of $\kk$-vector spaces.

\subsection{A finite reduction lemma}
\label{app:ss-finite-reduction}
We begin with a preliminary lemma needed for the argument below. 
Let $(\mathsf C^\bullet,d)$ be a bounded cochain complex of
finite-dimensional $\kk$-vector spaces.
Let $\mathcal Q$ be a finite indexing poset, and suppose that each term is equipped with a
decomposition
\[
\mathsf C^n
=
\bigoplus_{x\in\mathcal Q}\mathsf C_x^n 
\qquad
(n\in\mathbb Z).
\]

For $x,y\in\mathcal Q$, let
\[
d_{y,x}^n\colon
\mathsf C_x^n
\longrightarrow
\mathsf C_y^{n+1}
\]
denote the $(y,x)$-component of $d^n$.  Assume that $d$ is upper
triangular with respect to the $\mathcal Q$-indexed decomposition:
\begin{equation}
\label{eq:ss-upper-triangular-complex}
d_{y,x}^n\neq0
\quad\Longrightarrow\quad
x\leq y
\qquad
(x,y\in\mathcal Q,\ n\in\mathbb Z).
\end{equation}
We call $d$ strictly upper triangular if
$d_{x,x}^n=0$ for all $x\in\mathcal Q$ and $n\in\mathbb Z$.

Put $d_x:=d_{x,x}$. 
By \eqref{eq:ss-upper-triangular-complex}, $(\mathsf C_x^\bullet,d_x)$ is a cochain complex.

Consider the graded vector space 
\begin{equation}
\label{eq:ss-reduced-poset-graded-space}
\widetilde{\mathsf H}^n(\mathsf C)
=
\bigoplus_{x\in\mathcal Q}
H^n(\mathsf C_x^\bullet,d_x)
\qquad
(n\in\mathbb Z).
\end{equation}
The following lemma equips it with a differential that preserves
the triangular structure.

\begin{lemma}
\label{lem:ss-upper-triangular-cohomology}
There is a differential $\delta$ on
$\widetilde{\mathsf H}^\bullet(\mathsf C)$ such that $\delta$ is strictly upper
triangular with respect to
\eqref{eq:ss-reduced-poset-graded-space} and
$\bigl(\widetilde{\mathsf H}^\bullet(\mathsf C),\delta\bigr)
\simeq
(\mathsf C^\bullet,d)$ by a chain-homotopy equivalence.
\end{lemma}



\begin{proof}
For each $x\in\mathcal Q$, set
$Z_x^n:=\ker d_x^n$, $B_x^n:=\operatorname{im} d_x^{n-1}$.
Choose a complement $H_x^n$ of $B_x^n$ in $Z_x^n$, and then a complement
$L_x^n$ of $Z_x^n$ in $\mathsf C_x^n$, so that
\[
Z_x^n=B_x^n\oplus H_x^n,
\qquad
\mathsf C_x^n=H_x^n\oplus B_x^n\oplus L_x^n.
\]
Then $H_x^n\cong H^n(\mathsf C_x^\bullet,d_x)$ and
the restriction $d_x\colon L_x^n\to B_x^{n+1}$ is an isomorphism.
Applying Gaussian elimination to these invertible diagonal blocks, we obtain
a chain-homotopy equivalent complex whose underlying graded space is
\eqref{eq:ss-reduced-poset-graded-space}.

Because the original differential is upper triangular with respect to
the $\mathcal Q$-indexed decomposition, Gaussian elimination of a block
indexed by $z$ can modify a component from the $x$-part to the $y$-part
only when $x\leq z\leq y$.  Hence upper triangularity is preserved.
The diagonal components vanish, since each diagonal complex has been
replaced by its cohomology.  
Therefore the resulting differential is strictly upper triangular with
respect to the decomposition
\eqref{eq:ss-reduced-poset-graded-space}.
\end{proof}


For a complex $(\mathsf C^\bullet,d)$ as above, fix a differential
$\widetilde d_{\mathsf C}$ obtained by the construction in the proof of
\Cref{lem:ss-upper-triangular-cohomology}, and write
\begin{equation}
\label{eq:ss-upper-triangular-reduction-notation}
\bigl(
\widetilde{\mathsf H}^\bullet(\mathsf C),
\widetilde d_{\mathsf C}
\bigr)
\end{equation}
for the resulting cochain complex.

Next, we consider the interval subquotients associated with the
$\mathcal Q$-indexed decomposition.  For $x\leq y$ in $\mathcal Q$,
define the cochain complex
\begin{equation}
\label{eq:ss-poset-interval-complex}
\mathsf C^\bullet[x,y]
=
\bigoplus_{x\leq z}\mathsf C_z^\bullet
\bigg/
\bigoplus_{\substack{x\leq z\\ z\not\leq y}}
\mathsf C_z^\bullet
\end{equation}
as a quotient. 
Indeed, by \eqref{eq:ss-upper-triangular-complex}, the first direct
sum on the right is a subcomplex of $\mathsf C^\bullet$, and the
second is a subcomplex of the first.  We write $d^{[x,y]}$ for the
differential induced by $d$.  As a graded vector space,
$\mathsf C^\bullet[x,y]$ has the canonical decomposition
\begin{equation}
\label{eq:ss-poset-interval-decomposition}
\mathsf C^n[x,y]
\cong
\bigoplus_{x\leq z\leq y}\mathsf C_z^n
\qquad
(n\in\mathbb Z),
\end{equation}
and $d^{[x,y]}$ is upper triangular with respect to this
$[x,y]$-indexed decomposition.

Define
$\bigl(\widetilde{\mathsf H}^\bullet(\mathsf C)[x,y],
\widetilde d_{\mathsf C}^{[x,y]}\bigr)$
by applying the same quotient construction as in
\eqref{eq:ss-poset-interval-complex} to
$\bigl(\widetilde{\mathsf H}^\bullet(\mathsf C),
\widetilde d_{\mathsf C}\bigr)$.
Its underlying graded vector space is
\begin{equation}
\label{eq:ss-reduced-interval-complex}
\widetilde{\mathsf H}^n(\mathsf C)[x,y]
\cong
\bigoplus_{x\leq z\leq y}
H^n(\mathsf C_z^\bullet,d_z).
\end{equation}
The Gaussian eliminations in the proof of
\Cref{lem:ss-upper-triangular-cohomology} descend to these
subquotients by the same triangularity argument. Hence there is a
chain-homotopy equivalence
\begin{equation}
\label{eq:ss-reduced-interval-equivalence}
\bigl(
\widetilde{\mathsf H}^\bullet(\mathsf C)[x,y],
\widetilde d_{\mathsf C}^{[x,y]}
\bigr)
\simeq
\bigl(
\mathsf C^\bullet[x,y],d^{[x,y]}
\bigr).
\end{equation}
By construction,
\begin{equation}
\label{eq:ss-reduced-interval-locality}
(\widetilde d_{\mathsf C}^{[x,y]})_{y,x}
=
(\widetilde d_{\mathsf C})_{y,x}.
\end{equation}



\subsection{Reduction of the simple--simple derived Hom complex}
\label{app:ss-reduction-of-simple-simple-derived-hom-complex}
Let $P$ be a finite connected poset. 
Fix $S,C\in\Int(P)$.  By \Cref{prop:relative-auslander},
$\Lambda_P$ has finite global dimension.  We may therefore choose a
bounded projective resolution
$\mathcal P^\bullet(S)\to\LL(S)$ and a bounded injective coresolution
$\LL(C)\to\mathcal I^\bullet(C)$.

For a finite-dimensional $\Lambda_P$--$\Lambda_P$-bimodule $B$, we
consider the bounded cochain complex of finite-dimensional
$\kk$-vector spaces
\begin{equation}
\label{eq:ss-concrete-derived-evaluation}
\mathsf C^\bullet(B;S,C)
=
\operatorname{Tot}\Hom_{\Lambda_P}
\left(
B\otimes_{\Lambda_P}\mathcal P^\bullet(S),
\mathcal I^\bullet(C)
\right).
\end{equation}
The chosen resolutions give an isomorphism
\begin{equation}
\label{eq:ss-concrete-derived-evaluation-derived}
\mathsf C^\bullet(B;S,C)
\cong
\RHom_{\Lambda_P}
\left(
B\overset{\mathbf L}{\otimes}_{\Lambda_P}\LL(S),
\LL(C)
\right)
\quad\text{in }D(\kk).
\end{equation}
For later use, we record the matrix blocks of
\eqref{eq:ss-concrete-derived-evaluation}.  If $\PP(X)$ and
$\II(Y)$ occur as summands of the chosen resolutions, then
\begin{equation}
\label{eq:ss-bimodule-derived-evaluation-matrix-block}
\Hom_{\Lambda_P}
\left(
B\otimes_{\Lambda_P}\PP(X),
\II(Y)
\right)
\cong
\D(e_YBe_X).
\end{equation}
Under these identifications, the matrix components of the
differential are dual to the corresponding left and right
$\Lambda_P$-actions on $B$.


We first specialize to the regular bimodule $B=\Lambda_P$.  Then
\begin{equation}
\label{eq:ss-regular-derived-evaluation}
\mathsf C^\bullet(\Lambda_P;S,C)
\cong
\operatorname{Tot}\Hom_{\Lambda_P}
\left(
\mathcal P^\bullet(S),
\mathcal I^\bullet(C)
\right)
\end{equation}
and, by \eqref{eq:ss-concrete-derived-evaluation-derived}, it is
isomorphic to $\RHom_{\Lambda_P}(\LL(S),\LL(C))$ in $D(\kk)$.
By the definition of the total Hom complex, each term is given by
\begin{equation}
\label{eq:ss-total-hom-term}
\mathsf C^n(\Lambda_P;S,C)
=
\bigoplus_{q-p=n}
\Hom_{\Lambda_P}
\bigl(
\mathcal P^p(S),
\mathcal I^q(C)
\bigr).
\end{equation}
For the matrix blocks,
\eqref{eq:ss-bimodule-derived-evaluation-matrix-block} gives
\begin{equation}
\label{eq:ss-derived-evaluation-matrix-block}
\Hom_{\Lambda_P}\bigl(\PP(X),\II(Y)\bigr)
\cong
\D\bigl(e_Y\Lambda_Pe_X\bigr),
\qquad
X,Y\in\Int(P).
\end{equation}
Thus the admissible-component basis of $e_Y\Lambda_Pe_X$ gives a dual
basis of each matrix block.  For $T\in\Int(P)$, let
\begin{equation}
\label{eq:ss-support-component}
\mathsf C_T^n
=
\operatorname{span}_{\kk}
\left\{
\bigl((\rho_{Y,X}^{T})^{\op}\bigr)^\vee
\middle| 
\begin{array}{l}
T\in\Adm(Y,X),\\
\Hom_{\Lambda_P}(\PP(X),\II(Y))
\text{ is a matrix block in total degree }n
\end{array}
\right\}.
\end{equation}
We then
have
\begin{equation}
\label{eq:ss-support-decomposition}
\mathsf C^n(\Lambda_P;S,C)
=
\bigoplus_{T\in\Int(P)}
\mathsf C_T^n.
\end{equation}

Let $d$ denote the differential of
$\mathsf C^\bullet(\Lambda_P;S,C)$.  We claim that $d$ is upper
triangular with respect to the $\Int(P)$-indexed decomposition
\eqref{eq:ss-support-decomposition}.  Indeed, on each matrix block
the differential is dual to left or right multiplication in
$\Lambda_P$.  By \Cref{prop:admissible-basis-composition},
multiplication can only decrease the support, so dualizing gives
\begin{equation}
\label{eq:ss-support-upper-triangular}
d_{R,T}^n\neq0
\quad\Longrightarrow\quad
T\subseteq R
\end{equation}
as required. 

We may therefore apply
\Cref{lem:ss-upper-triangular-cohomology} to
$\mathsf C^\bullet(\Lambda_P;S,C)$.  
We write
\begin{equation}
\label{eq:ss-reduced-simple-simple-complex}
\bigl(
\widetilde{\mathsf H}^\bullet(S,C),
\widetilde d
\bigr)
:=
\bigl(
\widetilde{\mathsf H}^\bullet
    (\mathsf C^\bullet(\Lambda_P;S,C)),
\widetilde d_{\mathsf C(\Lambda_P;S,C)}
\bigr)
\end{equation}
for the resulting reduced complex.
By \Cref{lem:ss-upper-triangular-cohomology}, this complex is
chain-homotopy equivalent to
$\mathsf C^\bullet(\Lambda_P;S,C)$, and hence
\begin{equation}
\label{eq:ss-reduced-simple-simple-derived-hom}
\bigl(
\widetilde{\mathsf H}^\bullet(S,C),
\widetilde d
\bigr)
\cong
\RHom_{\Lambda_P}\bigl(\LL(S),\LL(C)\bigr)
\quad\text{in }D(\kk).
\end{equation}

We next describe the terms of
$\widetilde{\mathsf H}^\bullet(S,C)$ and the possible nonzero
components of its differential.

\begin{proposition}
\label{prop:ss-exact-support-reduced-terms}
For every $p\in\mathbb Z$, there is an isomorphism of
$\kk$-vector spaces
\begin{equation}
\label{eq:ss-exact-support-reduced-terms}
\widetilde{\mathsf H}^p(S,C)
\cong
\bigoplus_{\substack{T\in\mathcal W(S,C)\\
\omega(S,T,C)=p}}
\det(\partial_T^+S)
\otimes_\kk
\det(\partial_T^-C).
\end{equation}
\end{proposition}

\begin{proof}
We recall that 
\[
\widetilde{\mathsf H}^p(S,C)
=
\bigoplus_{T\in\Int(P)}
H^p(\mathsf C_T^\bullet,d_T).
\]

Fix $T\in\Int(P)$ and put $t=|T|$.
Set $B_T:=\DD(T)\otimes_\kk\D\NN(T)$.
By \Cref{prop:cardinality-bimodule-strata},
the isomorphism $\theta_t$ identifies $B_T$ with the subspace of
$J_t/J_{t-1}$ spanned by
$(\rho_{R,Q}^{T})^{\op}+J_{t-1}$ with
$T\in\Adm(R,Q)$ and $R,Q\in\Int(P)$.
Thus, for $X,Y\in\Int(P)$,
\eqref{eq:cardinality-bimodule-basis-correspondence} identifies
$e_YB_Te_X$ with the one-dimensional space spanned by
$(\rho_{Y,X}^{T})^{\op}+J_{t-1}$ when $T\in\Adm(Y,X)$, and both
spaces are zero otherwise.  Hence
\eqref{eq:ss-bimodule-derived-evaluation-matrix-block} and
\eqref{eq:ss-support-component} identify
$\mathsf C^n(B_T;S,C)$ with $\mathsf C_T^n$ for every $n$.

Since $\theta_t$ is a bimodule isomorphism, under these degreewise
identifications the left and right actions on $B_T$ are induced by
multiplication in $J_t/J_{t-1}$.  By
\Cref{prop:admissible-basis-composition}, every term whose
admissible-component support is a proper subset of $T$ lies in
$J_{t-1}$.  
Thus the differential on
$\mathsf C^\bullet(B_T;S,C)$ corresponds to $d_T$ under these
degreewise identifications. 
With $B_T=\DD(T)\otimes_\kk\D\NN(T)$, we obtain an isomorphism of cochain complexes
\[
(\mathsf C_T^\bullet,d_T)
\cong
\mathsf C^\bullet
\bigl(
\DD(T)\otimes_\kk\D\NN(T);S,C
\bigr).
\]

By \eqref{eq:ss-concrete-derived-evaluation-derived}, tensor--Hom
adjunction, and finite-dimensional $\kk$-duality, we have
\begin{equation}
\label{eq:ss-exact-support-factorization}
H^p(\mathsf C_T^\bullet,d_T)
\cong
\bigoplus_{i+j=p}
\Ext_{\Lambda_P}^i\bigl(\LL(S),\NN(T)\bigr)
\otimes_\kk
\Ext_{\Lambda_P}^j\bigl(\DD(T),\LL(C)\bigr).
\end{equation}
Applying \eqref{eq:simple-costandard-Ext} and
\eqref{eq:standard-simple-Ext} to
\eqref{eq:ss-exact-support-factorization} gives
\begin{equation}
\label{eq:ss-exact-support-diagonal-cohomology}
H^p(\mathsf C_T^\bullet,d_T)
\cong
\begin{cases}
\det(\partial_T^+S)\otimes_\kk\det(\partial_T^-C),
&
T\in\mathcal W(S,C)
\qand
p=\omega(S,T,C),
\\
0,
&
\text{otherwise}.
\end{cases}
\end{equation}
Substituting this into the preceding direct sum gives
\eqref{eq:ss-exact-support-reduced-terms}.
\end{proof}


By \eqref{eq:ss-exact-support-reduced-terms}, we identify the
underlying graded $\kk$-vector spaces of
$\widetilde{\mathsf H}^\bullet(S,C)$ and
$\mathsf K^\bullet(S,C)$.
It remains to identify the differential
$\widetilde d$.  Since $\widetilde d$ is strictly upper triangular,
a nonzero component $\widetilde d_{R,T}$ can occur only when
$T\subsetneq R$ and
$\omega(S,R,C)=\omega(S,T,C)+1$.  Together with
\eqref{eq:omega-boolean-rank}, this gives
\begin{equation}
\label{eq:ss-exact-support-cover-condition}
\widetilde d_{R,T}\neq0
\quad\Longrightarrow\quad
T\lessdot R
\text{ in }\mathcal W(S,C).
\end{equation}


\subsection{Cover components}
\label{app:ss-cover-components}

By \eqref{eq:ss-exact-support-cover-condition}, it remains to determine the cover components of $\widetilde d$.
We first establish a restriction lemma for standard modules and then
analyze the bimodule subquotient associated with a cover in
$\mathcal W(S,C)$.


\begin{lemma}
\label{lem:ss-standard-restriction-isomorphism}
Let $C,T,R\in\Int(P)$ with $T\in\cD(R)$.
If $(T,C),(R,C)\in \Upsilon^-(P)$ and
$\partial_T^-C=\partial_R^-C$, then the morphism $\varphi_{R,T}$ in \eqref{eq:standard-Hom-generator} is defined, and the map
\begin{equation}
\label{eq:ss-standard-restriction-map}
\Ext_{\Lambda_P}^p\bigl(\varphi_{R,T},\LL(C)\bigr)
\colon
\Ext_{\Lambda_P}^p\bigl(\DD(T),\LL(C)\bigr)
\longrightarrow
\Ext_{\Lambda_P}^p\bigl(\DD(R),\LL(C)\bigr)
\end{equation}
is an isomorphism for every $p\geq0$.
\end{lemma}

\begin{proof}
We first verify that $\varphi_{R,T}$ is defined.
We have $\Min(R)\subseteq T$.  Indeed,
$\Min(R)\cap C\subseteq C\subseteq T$, while
\[
\Min(R)\setminus C
\subseteq
\Min(R\setminus C)
=
\partial_R^-C
=
\partial_T^-C
\subseteq T.
\]
Since $T\in\cD(R)$, it follows that
$\Min(R)=\Min(T)$, as required. 

To prove the assertion, we use the following bar construction of a resolution of $\DD(V)$ for an interval $V\in\Int(P)$.
Put
$\mathsf B_0(V)=\PP(V)$ and, for $q\geq1$, let
\begin{equation}
\label{eq:ss-standard-bar-terms}
\mathsf B_q(V)
=
\bigoplus_{\substack{
U_0\subsetneq\cdots\subsetneq U_{q-1}\subsetneq V\\
U_i\in\cU(V)}}
\PP(U_0).
\end{equation}
The differential
$\mathsf B_q(V)\to\mathsf B_{q-1}(V)$ is defined as the alternating sum of the
following maps.  For $1\leq i\leq q-1$, deleting $U_i$ from
\[
U_0\subsetneq\cdots\subsetneq U_{q-1}\subsetneq V
\]
defines a map to the summand indexed by the resulting chain, acting as
the identity on $\PP(U_0)$.  Deleting $U_0$ uses the morphism
$\PP(U_0)\to\PP(U_1)$ induced by the inclusion $U_0\subsetneq U_1$.
When $q=1$, this last map is replaced by the morphism
$\PP(U_0)\to\PP(V)$ induced by $U_0\subsetneq V$.
Together with the canonical augmentation
$\mathsf B_0(V)=\PP(V)\to\DD(V)$, these maps form a chain complex
\begin{equation}
\label{eq:ss-standard-bar-resolution}
\cdots
\longrightarrow
\mathsf B_2(V)
\longrightarrow
\mathsf B_1(V)
\longrightarrow
\mathsf B_0(V)
\longrightarrow
\DD(V)
\longrightarrow
0.
\end{equation}

We claim that \eqref{eq:ss-standard-bar-resolution} is a projective
resolution of $\DD(V)$.  Every term is projective.  
After applying a primitive idempotent $e_X$ and fixing
$D\in\Adm(X,V)$, the support-$D$ part of $e_X\mathsf B_q(V)$ is
spanned by the chains
\[
D\subseteq U_0\subsetneq\cdots\subsetneq U_{q-1}\subsetneq V,
\qquad
U_i\in\cU(V).
\]
Thus these summands form the augmented simplicial chain complex of the
order complex of the half-open interval $[D,V)_{\cU(V)}$.
For $D\neq V$, this interval has the minimum $D$, so its order complex
is contractible and the augmented chain complex is exact.  For $D=V$,
there are no positive-degree terms and the augmentation identifies
the remaining basis vector with its image in $\DD(V)$.  Hence
\eqref{eq:ss-standard-bar-resolution} is exact.


Now let $C,T,R\in\Int(P)$ satisfy the hypotheses of the lemma. We
construct a chain map
\begin{equation}
\mathsf B_\bullet(R)\longrightarrow\mathsf B_\bullet(T)
\label{eq:ss-standard-bar-chain-map}
\end{equation}
lifting $\varphi_{R,T}$. In degree zero, we use the map
$(q_{R,T})_*\colon\PP(R)\to\PP(T)$. For $q\geq1$ and a chain
\[
U_0\subsetneq\cdots\subsetneq U_{q-1}\subsetneq R,
\]
let $V_0$ run through the connected components of $U_0\cap T$, and
for each $i>0$ let $V_i$ be the connected component of $U_i\cap T$
containing $V_0$.  The corresponding summand $\PP(U_0)$ is mapped to
$\PP(V_0)$ by the morphism induced by $q_{U_0,V_0}$, and we sum over
the possible $V_0$.  A term is taken to be zero if two consecutive
$V_i$ coincide.  
For every connected component $W$ of $U_1\cap T$, where for $q=1$ we read $U_1=R$ and $W=T$,
\Cref{prop:admissible-basis-composition} gives
\begin{equation}
\label{eq:ss-standard-bar-chain-map-composition}
q_{U_1,W}\circ j_{U_0,U_1}
=
\sum_{V\in\piZero(U_0\cap W)}
j_{V,W}\circ q_{U_0,V}.
\end{equation}
Thus the differential components corresponding to the deletion of
$U_0$ commute.  Compatibility with the remaining components follows
directly from deletion, so this defines the required chain map in \eqref{eq:ss-standard-bar-chain-map}.

For $V\in\{T,R\}$ and $q\geq1$, after applying
$\Hom_{\Lambda_P}(-,\LL(C))$ to $\mathsf B_q(V)$, a direct
summand $\PP(U_0)$ contributes only when $U_0=C$.
Thus the surviving
summands are indexed by chains
\[
C=U_0\subsetneq U_1\subsetneq\cdots\subsetneq U_{q-1}\subsetneq V.
\]
After omitting the fixed initial term $C$, these chains are precisely
the simplices of the order complex of the open interval
$(C,V)_{\cU(V)}$.

Applying $\Hom_{\Lambda_P}(-,\LL(C))$ to
\eqref{eq:ss-standard-bar-chain-map} gives a cochain map
\[
\Hom_{\Lambda_P}(\mathsf B_\bullet(T),\LL(C))
\longrightarrow
\Hom_{\Lambda_P}(\mathsf B_\bullet(R),\LL(C)).
\]
Under the identifications above, we claim that this is the simplicial pullback
induced by the order-preserving map
\begin{equation}
\label{eq:ss-standard-restriction-poset-map}
f\colon
(C,R)_{\cU(R)}
\longrightarrow
(C,T)_{\cU(T)},
\end{equation}
sending $U$ to the connected component of $U\cap T$ containing $C$.
We first check that the map $f$ is well-defined. 
Let $U\in(C,R)_{\cU(R)}$.  Since $U$ is a relative upset of $R$, $U\cap T$ is a relative upset of $T$.  Hence its connected component
$f(U)$ containing $C$ belongs to $\cU(T)$. 
It remains to show that
$C\subsetneq f(U)\subsetneq T$.
Since $U$ is connected and properly contains the relative upset $C$,
there exist $x\in U\setminus C$ and $c\in C$ with $x<c$.
Since $c\in T$ and $T\in\cD(R)$, we have $x\in T$.
Thus $C\subsetneq f(U)$.
To see that $f(U)\neq T$, suppose that $f(U)=T$.
Then $\Min(R)\subseteq T\subseteq U$.
Since $U$ is a relative upset of $R$, this gives $U=R$, a
contradiction. Thus, we have the map $f$ in \eqref{eq:ss-standard-restriction-poset-map}.
Moreover, $f$ is order preserving since $U\subseteq U'$ implies $f(U)\subseteq f(U')$.
Finally, it follows directly from the construction of
\eqref{eq:ss-standard-bar-chain-map} that the induced cochain map is the simplicial pullback associated with $f$.


On the other hand, we have an order-preserving map 
\begin{equation}
\label{eq:ss-standard-restriction-poset-inverse}
g\colon
(C,T)_{\cU(T)}
\longrightarrow
(C,R)_{\cU(R)},
\qquad
g(V)=V^{\uparrow_R}.
\end{equation}
Indeed, for every $V\in(C,T)_{\cU(T)}$, we have
$g(V)\in(C,R)_{\cU(R)}$ as follows.
First, $g(V)=V^{\uparrow_R}$ is a connected relative upset of $R$ and
$C\subsetneq g(V)$.
It remains to show that $g(V)\subsetneq R$.
Since $V\subsetneq T$ and
$\partial_T^-C=\Min(T\setminus C)$, there exists
$x\in\partial_T^-C$ with $x\notin V$.
By
$\partial_T^-C=\partial_R^-C=\Min(R\setminus C)$,
we have $x\in\Min(R\setminus C)$.
If $x\in V^{\uparrow_R}$, then some $v\in V$ satisfies $v\leq x$.
Since $C$ is a relative upset of $R$ and $x\notin C$, we have
$v\notin C$, and the minimality of $x$ in $R\setminus C$ gives
$v=x$, contradicting $x\notin V$.
Thus $g(V)\neq R$, and hence
$g(V)\in(C,R)_{\cU(R)}$.
Moreover, $g$ is order preserving, since
$V\subseteq V'$ implies
$V^{\uparrow_R}\subseteq (V')^{\uparrow_R}$.



Since $T\in\cD(R)$,
$(V^{\uparrow_R})\cap T=V$, while
$(f(U))^{\uparrow_R}\subseteq U$ because $U$ is a relative upset.
Hence
\begin{equation}
\label{eq:ss-standard-restriction-poset-homotopy}
f\circ g=\id,
\qquad
g\circ f\leq\id.
\end{equation}
Thus $f$ induces a homotopy equivalence of order complexes, and the
simplicial pullback induced by $f$ is a quasi-isomorphism.
Since the chain map
\eqref{eq:ss-standard-bar-chain-map} lifts $\varphi_{R,T}$,
\eqref{eq:ss-standard-restriction-map} is an isomorphism for $p>0$.

For $p=0$, the map is the identity if $T=R$.
If $T\subsetneq R$, then $C\neq T$ by the boundary equality, while
$C\neq R$ is immediate.
Hence both Hom groups vanish by the simple tops of the standard
modules.
\end{proof}



We next associate a bimodule subquotient with an inclusion
$T\subseteq R$ in $\Int(P)$.  Let
\[
B_{\leq R}
=
\operatorname{span}_\kk
\left\{
(\rho_{Q,U}^{W})^{\op}
\middle|
Q,U\in\Int(P),\ W\in\Adm(Q,U),\ W\subseteq R
\right\}
\]
and define
\begin{equation}
\label{eq:ss-cover-bimodule-subquotient}
B_{T,R}
=
B_{\leq R}
\bigg/
\operatorname{span}_\kk
\left\{
(\rho_{Q,U}^{W})^{\op}
\middle|
Q,U\in\Int(P),\ W\in\Adm(Q,U),\
W\subseteq R,\ T\nsubseteq W
\right\}.
\end{equation}
By \Cref{prop:admissible-basis-composition}, the numerator and
denominator in \eqref{eq:ss-cover-bimodule-subquotient} are
$\Lambda_P$--$\Lambda_P$-subbimodules.  
The induced admissible-component basis of $B_{T,R}$ therefore
consists of the basis vectors whose supports $W$ satisfy
$T\subseteq W\subseteq R$.


With the decompositions fixed above, the admissible-component basis
identifies
\begin{equation}
\label{eq:ss-cover-local-complex}
\mathsf C^\bullet(B_{T,R};S,C)
\cong
\mathsf C^\bullet(\Lambda_P;S,C)[T,R]
\end{equation}
as cochain complexes.  
Indeed, on each matrix block, both sides retain exactly the dual
admissible-component basis vectors with supports $W$ satisfying
$T\subseteq W\subseteq R$, together with the same differential
components between them.

\begin{proposition}
\label{prop:ss-cover-bimodule-acyclicity}
Let $S,C\in\Int(P)$ and let $T\lessdot R$ in
$\mathcal W(S,C)$.  Then the complex
$\mathsf C^\bullet(B_{T,R};S,C)$ is acyclic.
\end{proposition}


\begin{proof}
Since $T\lessdot R$ in $\mathcal W(S,C)$, either
\eqref{eq:ss-upper-cover-data} or
\eqref{eq:ss-lower-cover-data} holds.

Assume first that \eqref{eq:ss-upper-cover-data} holds.
We first note that $\Min(R)\subseteq T$.  Indeed,
$\Min(R)\cap C\subseteq T$, while
\[
\Min(R)\setminus C
\subseteq
\Min(R\setminus C)
=
\partial_R^-C
=
\partial_T^-C
\subseteq T.
\]
Since $T$ is convex, it follows that $T\in\cD(R)$.  Consequently,
every interval $W$ with $T\subseteq W\subseteq R$ also belongs to
$\cD(R)$.  Moreover, since $\Min(R)\subseteq T\subseteq W$, we have $\Min(W)=\Min(R)$. In particular, the map $\varphi_{R,W}$ is defined.


Put $M_{T,R}:=e_RB_{T,R}$, regarded as a right $\Lambda_P$-module.
Since $\Min(R)\subseteq T$, no proper relative upset of $R$ contains
$T$.  By \eqref{eq:standard-kernel-generators} and
\Cref{prop:admissible-basis-composition}, the left
$\Lambda_P$-action on $B_{T,R}$ gives
$\KK(R)M_{T,R}=0$ inside $B_{T,R}$.
Multiplication therefore induces a bimodule map
\[
\eta_{T,R}\colon
\DD(R)\otimes_\kk M_{T,R}
\longrightarrow
B_{T,R}.
\]


Filter $B_{T,R}$ by the cardinality of the admissible-component
support, and give $M_{T,R}=e_RB_{T,R}$ the induced filtration.
We use the resulting filtration on
$\DD(R)\otimes_\kk M_{T,R}$ through the second factor.
Since multiplication cannot increase 
the admissible-component support, $\eta_{T,R}$ preserves these filtrations.

By \Cref{prop:cardinality-bimodule-strata} and
\eqref{eq:cardinality-bimodule-basis-correspondence}, for
$T\subseteq W\subseteq R$ and $t=|W|$, the summands corresponding
to $W$ in the $t$th successive quotients of the source and target
of $\eta_{T,R}$ are identified with
$\DD(R)\otimes_\kk\D\NN(W)$ and
$\DD(W)\otimes_\kk\D\NN(W)$, respectively.
Under these identifications, the map 
$\DD(R)\otimes_\kk\D\NN(W)\to \DD(W)\otimes_\kk\D\NN(W)$ 
induced by $\eta_{T,R}$ is exactly
$\varphi_{R,W}\otimes_\kk\id_{\D\NN(W)}$.

We claim that applying $\mathsf C^\bullet(-;S,C)$ to this map gives
a quasi-isomorphism. Since $C\in\cU(R)$ and
$C\subseteq W\subseteq R$, we have $C\in\cU(W)$. Moreover,
\[
\partial_W^-C
=
\partial_R^-C
=
\partial_T^-C
=
\Min(W\setminus C).
\]
Indeed, $\partial_R^-C=\partial_T^-C\subseteq T\subseteq W$, and
the convexity of $W$ in $R$ gives
$\partial_W^-C=\partial_R^-C$. Since
$\partial_R^-C=\Min(R\setminus C)$ and $W\in\cD(R)$, this set is
also $\Min(W\setminus C)$. Thus $(W,C)\in\Upsilon^-(P)$.
For $W\neq R$, \Cref{lem:ss-standard-restriction-isomorphism} shows
that $\Ext_{\Lambda_P}^p(\varphi_{R,W},\LL(C))$ is an isomorphism
for every $p\geq0$, while for $W=R$ the map is the identity.
By \eqref{eq:ss-concrete-derived-evaluation-derived} and
tensor--Hom adjunction, the claim follows.

Thus, after applying $\mathsf C^\bullet(-;S,C)$, the map induced by
$\eta_{T,R}$ is a quasi-isomorphism on every successive quotient.
An induction on the finite filtrations shows that the map induced by
$\eta_{T,R}$ is a quasi-isomorphism
\begin{equation}
\label{eq:ss-cover-derived-evaluation-comparison}
\mathsf C^\bullet(B_{T,R};S,C)
\longrightarrow
\mathsf C^\bullet(\DD(R)\otimes_\kk M_{T,R};S,C).
\end{equation}


It remains to show that the complex on the right is acyclic.
Since the right $\Lambda_P$-action on
$\DD(R)\otimes_\kk M_{T,R}$ is carried by the second factor, we have
\[
\bigl(\DD(R)\otimes_\kk M_{T,R}\bigr)
\overset{\mathbf L}{\otimes}_{\Lambda_P}\LL(S)
\cong
\DD(R)\otimes_\kk
\left(
M_{T,R}
\overset{\mathbf L}{\otimes}_{\Lambda_P}\LL(S)
\right).
\]
By \eqref{eq:ss-concrete-derived-evaluation-derived}, it remains to prove
\begin{equation}
\label{eq:ss-cover-one-sided-vanishing}
M_{T,R}
\overset{\mathbf L}{\otimes}_{\Lambda_P}\LL(S)
\cong0
\quad\text{in }D(\kk).
\end{equation}



To prove \eqref{eq:ss-cover-one-sided-vanishing}, we construct a projective resolution of the right module $M_{T,R}$.
Put $\mathsf F_0=e_R\Lambda_P$ and, for $q\geq1$, let
\[
\mathsf F_q
=
\bigoplus_{\substack{
D_0\subsetneq\cdots\subsetneq D_{q-1}\\
D_i\in\cD(R),\ T\nsubseteq D_i}}
e_{D_0}\Lambda_P.
\]
We define the differential
$\mathsf F_q\to\mathsf F_{q-1}$ as the alternating sum of the
following maps.  Deleting $D_i$ for $1\leq i\leq q-1$ acts as the
identity on $e_{D_0}\Lambda_P$, while deleting $D_0$ is given by left
multiplication by $q_{D_1,D_0}^{\op}$, where $D_1=R$ when $q=1$.
Together with the quotient
$e_R\Lambda_P\twoheadrightarrow M_{T,R}$, these maps form an
augmented complex of projective right modules.

We verify its exactness on the admissible-component basis.
After multiplying on the right by a primitive idempotent $e_X$, fix
a support $W\subseteq R$ occurring in the admissible-component basis.
If $T\subseteq W$, no support-$W$ basis vector occurs in
$\mathsf F_q$ for $q\geq1$, since such a vector would require
$W\subseteq D_0$ and hence $T\subseteq D_0$, contrary to the
indexing condition. 
In degree zero, the augmentation maps the support-$W$ basis vector of $e_R\Lambda_P$ to its class in $M_{T,R}$, and this class is nonzero.
If $T\nsubseteq W$, the terms with support $W$ form the augmented
simplicial chain complex of the order complex of
\[
\left\{
D\in\cD(R)
\mid
W\subseteq D,\ T\nsubseteq D
\right\}.
\]
The admissibility conditions imply $W\in\cD(R)$, so this poset has
minimum $W$ and hence has contractible order complex.
Thus $\mathsf F_\bullet\to M_{T,R}$ is a projective resolution.
Therefore
\begin{equation}
\label{eq:ss-cover-derived-tensor-resolution}
M_{T,R}
\overset{\mathbf L}{\otimes}_{\Lambda_P}\LL(S)
\cong
\mathsf F_\bullet\otimes_{\Lambda_P}\LL(S)
\quad\text{in }D(\kk).
\end{equation}
We now show that the complex on the right is acyclic. 
For every $D\in\Int(P)$, we have
\[
e_D\Lambda_P\otimes_{\Lambda_P}\LL(S)
\cong
e_D\LL(S)
=
\begin{cases}
\kk, & D=S,\\
0, & D\neq S.
\end{cases}
\]
Thus, in positive degrees, the only summands of
$\mathsf F_\bullet\otimes_{\Lambda_P}\LL(S)$ that survive are those
indexed by chains with $D_0=S$.

In degree zero,
\[
\mathsf F_0\otimes_{\Lambda_P}\LL(S)
=
e_R\Lambda_P\otimes_{\Lambda_P}\LL(S)
\cong e_R\LL(S)=0,
\]
since $S\subseteq T\subsetneq R$.

If $S=T$, the resulting complex is zero.  Indeed, there are no
positive-degree summands with $D_0=S$, since the indexing condition
$T\nsubseteq D_0$ would then fail, while the degree-zero term vanishes
as above.

Suppose that $S\subsetneq T$.  For $q\geq1$, the surviving summands in
$\mathsf F_q\otimes_{\Lambda_P}\LL(S)$ are indexed by chains
\[
S=D_0\subsetneq D_1\subsetneq\cdots\subsetneq D_{q-1},
\qquad
D_i\in\cD(R),\quad T\nsubseteq D_i.
\]
After omitting the fixed initial term $D_0=S$, these are precisely the
chains in the poset
\[
\cD_{S,T}(R)
:=
\left\{
D\in\cD(R)
\mid
S\subsetneq D,\ T\nsubseteq D
\right\}.
\]
The induced differential is the alternating deletion differential on
these chains.  
Thus, up to a degree shift, the remaining complex is the augmented
simplicial chain complex of the order complex of $\cD_{S,T}(R)$. 
To prove the acyclicity, it therefore suffices to show that
$\cD_{S,T}(R)$ is contractible.
Let $a$ be the element in \eqref{eq:ss-upper-cover-data} and set  
\[
c(D):=(D\cup\{a\})^{\downarrow_R}
\qquad
(D\in\cD_{S,T}(R)).
\]
This is a connected relative downset of $R$ containing $D$.
Since $T\nsubseteq D$ and
$\partial_T^+S=\Max(T\setminus S)$, there is
$b\in\partial_T^+S$ with $b\notin D$. 
Since
$a\in\partial_R^+S\setminus\partial_T^+S$
and $b\in\partial_T^+S\subseteq\partial_R^+S$, we have
$a\neq b$ and
$a,b\in\partial_R^+S=\Max(R\setminus S)$.
In particular, $b\nleq a$.
Moreover, since $D$ is a relative downset of $R$ and $b\notin D$,
we have $b\notin(D\cup\{a\})^{\downarrow_R}=c(D)$.
Thus $T\nsubseteq c(D)$, so $c(D)\in\cD_{S,T}(R)$.
The map $c$ is an ascending closure operator on $\cD_{S,T}(R)$.
Its image has the least element $(S\cup\{a\})^{\downarrow_R}$, and
hence $\cD_{S,T}(R)$ is contractible.
Therefore the complex above is acyclic, proving
\eqref{eq:ss-cover-one-sided-vanishing}.

Finally, the case \eqref{eq:ss-lower-cover-data} follows from the
preceding one by applying order reversal $P\mapsto P^{\op}$, with
$S$ and $C$ interchanged, and then applying $\kk$-duality, which
preserves acyclicity.
\end{proof}



\begin{proposition}
\label{prop:ss-cover-component}
Let $S,C\in\Int(P)$ and let $T\lessdot R$ in
$\mathcal W(S,C)$.  Then the component
\[
\widetilde d_{R,T}:
\det(\partial_T^+S)\otimes_\kk\det(\partial_T^-C)
\longrightarrow
\det(\partial_R^+S)\otimes_\kk\det(\partial_R^-C)
\]
is an isomorphism.
\end{proposition}

\begin{proof}
By \eqref{eq:ss-cover-local-complex} and
\eqref{eq:ss-reduced-interval-equivalence},
$\widetilde{\mathsf H}^\bullet(S,C)[T,R]$ is chain-homotopy
equivalent to
$\mathsf C^\bullet(B_{T,R};S,C)$.
By \eqref{eq:ss-exact-support-diagonal-cohomology}, the nonzero
summands of $\widetilde{\mathsf H}^\bullet(S,C)[T,R]$ are indexed by
the intervals $W\in\mathcal W(S,C)$ satisfying
$T\subseteq W\subseteq R$.
Since $T\lessdot R$ in $\mathcal W(S,C)$, these are precisely
$T$ and $R$.
By \eqref{eq:omega-boolean-rank}, their degrees differ by one.
Hence $\widetilde{\mathsf H}^\bullet(S,C)[T,R]$ is the two-term complex
\[
\det(\partial_T^+S)\otimes_\kk\det(\partial_T^-C)
\longrightarrow
\det(\partial_R^+S)\otimes_\kk\det(\partial_R^-C).
\]
Its differential is $\widetilde d_{R,T}$ by
\eqref{eq:ss-reduced-interval-locality}.
The complex $\mathsf C^\bullet(B_{T,R};S,C)$ is acyclic by
\Cref{prop:ss-cover-bimodule-acyclicity}, so this two-term complex is
acyclic. Therefore $\widetilde d_{R,T}$ is an isomorphism.
\end{proof}


We can now complete the proof of
\Cref{prop:simple-simple-combinatorial-model}.


\begin{proof}[Proof of \Cref{prop:simple-simple-combinatorial-model}]
By \eqref{eq:simple-simple-combinatorial-complex} and
\eqref{eq:ss-exact-support-reduced-terms}, we identify the underlying
graded $\kk$-vector spaces of $\mathsf K^\bullet(S,C)$ and
$\widetilde{\mathsf H}^\bullet(S,C)$. If
$\mathcal W(S,C)=\varnothing$, both complexes are zero, and the result
follows from \eqref{eq:ss-reduced-simple-simple-derived-hom}.
Suppose that $\mathcal W(S,C)$ is nonempty.

By \eqref{eq:ss-exact-support-cover-condition} and
\Cref{prop:ss-cover-component}, the nonzero components of
$\widetilde d$ are precisely those corresponding to covers
$T\lessdot R$ in $\mathcal W(S,C)$. For each such cover, let
$d^{\mathsf K}_{R,T}$ denote the component defined in
\eqref{eq:simple-simple-combinatorial-differential}.
Both $d^{\mathsf K}_{R,T}$ and $\widetilde d_{R,T}$ are isomorphisms
between the same one-dimensional vector spaces. Hence there is a
unique $\alpha_{R,T}\in\kk^\times$ such that
$\widetilde d_{R,T}=\alpha_{R,T}d^{\mathsf K}_{R,T}$.

By \Cref{prop:intermediate-boolean-structure},
$\mathcal W(S,C)$ is a Boolean lattice. Let $T_{\min}$ be its
minimum, put $c_{T_{\min}}=1$, and define $c_T$ as the product of the
scalars $\alpha$ along a saturated chain from $T_{\min}$ to $T$.
This is independent of the chosen chain. Indeed, any two saturated
chains are related by moves across squares in the Hasse diagram of
$\mathcal W(S,C)$, and the relations
$(d^{\mathsf K})^2=0$ and $\widetilde d^{\,2}=0$ show that the
products of the scalars $\alpha$ along the two paths around each
square agree. Thus $c_R=\alpha_{R,T}c_T$ for every cover
$T\lessdot R$.
Multiplication by $c_T$ on the summand indexed by $T$ therefore gives
an isomorphism of cochain complexes
\[
\bigl(\mathsf K^\bullet(S,C),d^{\mathsf K}\bigr)
\cong
\bigl(\widetilde{\mathsf H}^\bullet(S,C),\widetilde d\bigr).
\]
Together with \eqref{eq:ss-reduced-simple-simple-derived-hom}, this
gives
\[
\bigl(\mathsf K^\bullet(S,C),d^{\mathsf K}\bigr)
\cong
\RHom_{\Lambda_P}\bigl(\LL(S),\LL(C)\bigr)
\quad\text{in }D(\kk),
\]
as required.
\end{proof}

